\maketableofcontents

\chapter*{ВВЕДЕНИЕ}
\addcontentsline{toc}{chapter}{ВВЕДЕНИЕ}

Компилятор является одним из ключевых элементов системного программного обеспечения. Он связывает язык высокого уровня с конкретной аппаратной и программной платформой, выполняя лексический анализ, синтаксический разбор, построение внутреннего представления программы и генерацию машинно-зависимого кода. Разработка даже небольшого компилятора позволяет изучить устройство трансляторов, представление выражений и операторов, организацию таблицы символов, соглашения о вызовах функций и особенности целевой архитектуры.

В данной курсовой работе рассматривается язык B. Язык B был разработан Кеном Томпсоном и Деннисом Ритчи в Bell Labs как потомок BCPL и непосредственный предшественник языка C. В историческом варианте B является компактным бестиповым языком, предназначенным прежде всего для системного программирования. Его синтаксис включает внешние определения, функции, автоматические и внешние переменные, векторы, адресную арифметику, условные операторы, циклы, оператор выбора, переходы и развитую систему выражений.

Актуальность работы заключается в том, что язык B достаточно мал для полного описания в рамках курсового проекта, но при этом содержит основные конструкции императивного языка программирования: объявления, области видимости, функции, рекурсию, указатели, массивы, условные и циклические операторы. Поэтому реализация компилятора B позволяет последовательно рассмотреть основные этапы построения транслятора без чрезмерного усложнения грамматики.

Целью курсовой работы является разработка компилятора языка B на языке C, выполняющего трансляцию исходного кода в исполняемый машинный код для архитектуры x86-64 под операционную систему Linux.

Для достижения поставленной цели необходимо решить следующие задачи:

\begin{itemize}
  \item изучить справочное руководство по языку B и выделить основные синтаксические конструкции языка;
  \item описать грамматику поддерживаемого подмножества языка B;
  \item разработать лексический анализатор для выделения ключевых слов, идентификаторов, литералов, операторов и разделителей;
  \item разработать синтаксический анализатор для объявлений, выражений и операторов языка;
  \item реализовать таблицу символов для глобальных объектов, локальных переменных, параметров функций, внешних имен и меток;
  \item построить внутреннее представление программы в виде абстрактного синтаксического дерева;
  \item реализовать генерацию ассемблерного кода x86-64 для ОС Linux;
  \item обеспечить сборку исполняемого файла с использованием системного ассемблера и компоновщика;
  \item подготовить тестовые программы на языке B и проверить соответствие результатов ожидаемому выводу.
\end{itemize}

Пояснительная записка состоит из аналитической, конструкторской и технологической частей. В аналитической части приведено описание языка B и требований к компилятору. В конструкторской части рассмотрена архитектура транслятора и внутренние структуры данных. В технологической части описана реализация, приведены примеры программ, фрагменты полученного ассемблерного кода и результаты тестирования.
